\scanpage{78}
\margin{\tikz[baseline=-.3ex]{\draw (0,0) circle (.035) (0,.03)--(.30,.52);}} % slanted
$H^2(S,A)$ \deleted{\unsure{de} $H^2$} (dont le $T_\ell$ est $H^2(S,T_\ell(A))$) % 1
n'est pas néc nul, \deleted{\unsure{si}} \inserted{comme on voit} si $S$ est % 2
une courbe lisse projective géom connexe sur $\mathbb{F}_p$, % 3
et $A=J\times S$, où $J$ est la jacobienne de $S$. % 4
On a en tous cas % 5
\[
\begin{aligned}
 H^2(S,T_\ell(A)) &\simeq \unsure{\mathbb{H}|H}^{1}(\bar S,T_\ell(\bar A))_{\pi}
 =\deleted{T_\ell J(\overline{\mathbb{F}}_p)\otimes}
 \inserted{\bigl[T_\ell(\bar J)\otimes T_\ell(\bar J)\bigr]_{\pi}}\\ % 6
 &\simeq\operatorname{End}(J)\otimes_{\deleted{\unsure{\mathbb{Q}}}\mathbb{Z}}\mathbb{Q}_{\ell}
 \quad\deleted{\text{\unsure{qui}}}\bigl(\simeq
 \bigl(\text{\unsure{Algèbre|classes} de corr. }\inserted{\text{\unsure{divisorielles}}}\text{ sur la }C\bigr)
 \otimes_{\mathbb{Z}}\mathbb{Q}_{\ell} % 7
\end{aligned}
\]
qui est non nul ----. % 8
% archive sheet number: 77
